\SourceLeaf{407}%
\[
\text{V.}\quad .\ .\ .\ .\ .\ .\ .\qquad b=90^0,\quad A\gtrless90^0
\]
\par
1) \(c=90^0\). Then \(A\) is the pole of \(c\), so \(b=90^0\), and \(a=A\). The equation
\[
\cos a=\cos b.\cos c+\sin b.\sin c.\cos A
\]
is evident.

2) \(c<90^0\). Extend \(c\) beyond \(B\) until its length is \(90^0\). Considering the triangle then gives
\[
\cos a=\cos(90^0-c).\cos A+\sin(90^0-c).\sin A.\cos90^0
\]
or
\[
\cos a=\sin c.\cos A
\]
so the formula is also correct in this case.

3) \(c>90^0\). From \(c\), an arc of \(90^0\) is cut off at the point \(R\). Considering the triangle \(BRC\) then gives
\[
\cos a=\cos(c-90^0).\cos A+\sin(c-90^0).\sin A.\cos90^0
\]
or
\[
\cos a=\sin c.\cos A
\]
as in the preceding case.

The cases in which \(b<90^0\) or \(b>90^0\), with \(c=90^0\) at the same time, are essentially the same as the preceding two cases.

Listing all the possible cases, we obtain the proof as follows
\begin{center}
\begin{tabular}{ll}
Both sides less than \(90^0\)&from I\\
Both sides greater than \(90^0\)&from II\\
One greater and the other less than \(90^0\)&from III\\
Both \(=90^0\)&from IV and V.\,1.\\
One \(=90^0\) and the other less&from IV and V.\,2.\\
One \(=90^0\) and the other greater&from IV and V.\,3.\\
\end{tabular}
\end{center}
\hspace*{\parindent}We can therefore take the following formula to hold in general
\[
\text{(A)}\quad .\ .\ .\ .\ .\quad\cos a=\cos b.\cos c+\sin b.\sin c.\cos A
\]
and likewise
